Un haut-parleur placé dans une salle émet un son pur de \qty{\v{f}}{Hz} vers une porte ouverte de \qty{\v{larg}}{m} de largeur. Dans le corridor, loin de la porte, on mesure l'intensité sonore en fonction de l'angle par rapport à l'axe de la porte. La vitesse du son est de \qty{\v{v}}{m/s}.

\begin{enumerate}[label=\alph*)]
    \item Déterminez la longueur d'onde du son.
    \item À quel angle de l'axe de la porte l'intensité sonore est-elle nulle pour la première fois ?
    \item Combien de directions d'intensité nulle y a-t-il de chaque côté de l'axe ?
    \item Sous quelle fréquence n'y a-t-il plus aucune direction d'intensité nulle ?
\end{enumerate}
